% ============================================================================
%  Chapter 9 --- Cloud Deployment (HCS Mapping)
% ============================================================================
\chapter{Cloud Deployment Strategy}
\label{ch:cloud}

\section{Introduction}

A core strength of this project is its cloud portability. The platform was designed from inception to run on Docker Compose locally and deploy to Huawei Cloud Stack (HCS) with minimal modifications. This chapter describes the HCS mapping, the deployment architecture, and the configuration changes required for cloud deployment.

\section{Local-to-Cloud Mapping}

% ── Figure: HCS Mapping ──────────────────────────────────────────────────────
\begin{figure}[H]
\centering
\begin{tikzpicture}[node distance=1.5cm and 3.5cm]

  % Local side
  \node[font=\large\bfseries, text=darkbg] (local-title) at (0, 5) {Local (Docker Compose)};
  
  \node[service=ossblue, minimum width=3.5cm] (l-api) at (0, 3.5) {
    \texttt{api-gateway}\\{\scriptsize Docker container}
  };
  \node[service=aiviolet, minimum width=3.5cm] (l-ai) at (0, 2) {
    \texttt{ai-service}\\{\scriptsize Docker container}
  };
  \node[service=cemteal, minimum width=3.5cm] (l-worker) at (0, 0.5) {
    \texttt{pipeline-worker}\\{\scriptsize Docker run-once}
  };
  \node[datastore=datalake, minimum width=3.5cm, minimum height=1cm] (l-minio) at (0, -1.2) {
    \textbf{MinIO}\\{\scriptsize Object Storage}
  };
  \node[datastore=bssgreen, minimum width=3.5cm, minimum height=1cm] (l-pg) at (0, -2.8) {
    \textbf{PostgreSQL 16}\\{\scriptsize Database}
  };
  
  % HCS side  
  \node[font=\large\bfseries, text=huaweired!80!black] (hcs-title) at (8, 5) {Huawei Cloud Stack (HCS)};
  
  \node[service=ossblue, minimum width=3.5cm] (h-api) at (8, 3.5) {
    \textbf{ECS / CCE}\\{\scriptsize api-gateway}
  };
  \node[service=aiviolet, minimum width=3.5cm] (h-ai) at (8, 2) {
    \textbf{ECS / CCE}\\{\scriptsize ai-service}
  };
  \node[service=cemteal, minimum width=3.5cm] (h-worker) at (8, 0.5) {
    \textbf{ECS / CCE}\\{\scriptsize pipeline-worker}
  };
  \node[datastore=datalake, minimum width=3.5cm, minimum height=1cm] (h-obs) at (8, -1.2) {
    \textbf{OBS}\\{\scriptsize Object Storage Service}
  };
  \node[datastore=bssgreen, minimum width=3.5cm, minimum height=1cm] (h-rds) at (8, -2.8) {
    \textbf{RDS}\\{\scriptsize PostgreSQL Managed}
  };
  
  % HCS boundary
  \begin{scope}[on background layer]
    \node[fit=(hcs-title)(h-api)(h-ai)(h-worker)(h-obs)(h-rds),
          draw=huaweired!40, fill=huaweired!3, rounded corners=8pt,
          inner sep=12pt, line width=0.8pt] (hcsbox) {};
  \end{scope}
  
  % VPC label
  \node[below=0.2cm of hcsbox.south, font=\scriptsize\itshape, text=huaweired!60!black] {
    VPC (Virtual Private Cloud) --- network isolation
  };
  
  % Mapping arrows
  \draw[myarrow=gray, line width=1.2pt, dashed] (l-api.east) -- (h-api.west)
    node[arrlabel, midway, above] {same image};
  \draw[myarrow=gray, line width=1.2pt, dashed] (l-ai.east) -- (h-ai.west)
    node[arrlabel, midway, above] {same image};
  \draw[myarrow=gray, line width=1.2pt, dashed] (l-worker.east) -- (h-worker.west)
    node[arrlabel, midway, above] {same image};
  \draw[myarrow=datalake!70, line width=1.2pt] (l-minio.east) -- (h-obs.west)
    node[arrlabel, midway, above] {S3 API compatible};
  \draw[myarrow=bssgreen!70, line width=1.2pt] (l-pg.east) -- (h-rds.west)
    node[arrlabel, midway, above] {same protocol};
    
\end{tikzpicture}
\caption{Local Docker Compose to Huawei Cloud Stack mapping. Application containers run unchanged; only infrastructure endpoints change.}
\label{fig:hcs-mapping}
\end{figure}

\section{Mapping Details}

\begin{table}[H]
\centering
\caption{Detailed local-to-HCS component mapping.}
\label{tab:hcs-mapping}
\begin{tabularx}{\textwidth}{@{}lllX@{}}
\toprule
\textbf{Local Component} & \textbf{HCS Service} & \textbf{Change Required} & \textbf{Notes} \\
\midrule
Docker containers & ECS or CCE & Image registry URL & Same Docker images, pushed to HCS container registry \\
MinIO (:9000) & OBS & \texttt{S3\_ENDPOINT} env var & OBS is S3-compatible; boto3 client works unchanged \\
PostgreSQL 16 & RDS for PostgreSQL & \texttt{DATABASE\_URL} env var & HCS manages backups, HA, scaling \\
Docker volumes & OBS + RDS & Automatic & Managed services handle persistence \\
Docker network & VPC & VPC subnet config & Network isolation and security groups \\
Port exposure & ELB & Load balancer config & Elastic Load Balancer for API Gateway \\
\bottomrule
\end{tabularx}
\end{table}

\section{Configuration Changes}

The transition from local to HCS requires exactly \textbf{three environment variable changes}:

\begin{lstlisting}[language=bash, caption={Environment changes for HCS deployment.}]
# Local configuration
DATABASE_URL=postgresql://telecom:pw@postgres:5432/telecom_intel
S3_ENDPOINT=http://minio:9000
S3_ACCESS_KEY=minio

# HCS configuration
DATABASE_URL=postgresql://telecom:pw@rds-xxx.huaweicloud.com:5432/telecom_intel
S3_ENDPOINT=https://obs.xx-region.huaweicloud.com
S3_ACCESS_KEY=<hcs-ak>
\end{lstlisting}

No code changes are required. The same Docker images run on both environments.

\section{HCS Deployment Architecture}

% ── Figure: HCS Target Architecture ──────────────────────────────────────────
\begin{figure}[H]
\centering
\begin{tikzpicture}[node distance=1.2cm and 2cm]

  % VPC boundary
  \node[font=\footnotesize\bfseries, text=huaweired!70!black] (vpclabel) at (4.5, 5.5) {Huawei Cloud Stack --- VPC};
  
  % ELB
  \node[service=cvmorange, minimum width=8cm, minimum height=0.8cm] (elb) at (4.5, 4.5) {
    \textbf{ELB} (Elastic Load Balancer) --- HTTPS termination
  };
  
  % Application subnet
  \node[service=ossblue, minimum width=2.5cm] (api) at (1.5, 3) {
    \textbf{ECS}\\api-gateway
  };
  \node[service=aiviolet, minimum width=2.5cm] (ai) at (4.5, 3) {
    \textbf{ECS}\\ai-service
  };
  \node[service=cemteal, minimum width=2.5cm] (worker) at (7.5, 3) {
    \textbf{ECS}\\pipeline-worker
  };
  
  % Data subnet
  \node[datastore=datalake, minimum width=3cm, minimum height=1cm] (obs) at (2.5, 0.8) {
    \textbf{OBS}\\raw / proc / curated
  };
  \node[datastore=bssgreen, minimum width=3cm, minimum height=1cm] (rds) at (6.5, 0.8) {
    \textbf{RDS}\\PostgreSQL (managed)
  };
  
  % VPC box
  \begin{scope}[on background layer]
    \node[fit=(vpclabel)(elb)(api)(ai)(worker)(obs)(rds),
          draw=huaweired!40, fill=huaweired!2, rounded corners=10pt,
          inner sep=15pt, line width=1pt] {};
  \end{scope}
  
  % Arrows
  \draw[myarrow] (elb) -- (api);
  \draw[myarrow=aiviolet!60] (worker) -- (ai);
  \draw[myarrow=cemteal!60] (worker) -- (obs);
  \draw[myarrow=cemteal!60] (worker) -- (rds);
  \draw[myarrow=ossblue!60] (api) -- (rds);
  \draw[myarrow=aiviolet!60] (ai) -- (obs);
  
  % External user
  \node[external, above=0.8cm of elb] (ext) {CVM / Analyst\\{\scriptsize HTTPS}};
  \draw[myarrow=cvmorange!70, line width=1.2pt] (ext) -- (elb);
  
\end{tikzpicture}
\caption{Target HCS deployment architecture with VPC, ELB, ECS, OBS, and RDS.}
\label{fig:hcs-target}
\end{figure}

\section{Deployment Steps}

\begin{enumerate}
  \item \textbf{Create VPC} with two subnets: application (for ECS instances) and data (for RDS/OBS).
  \item \textbf{Provision RDS} (PostgreSQL 16, primary + standby for HA).
  \item \textbf{Create OBS buckets}: \texttt{raw}, \texttt{processed}, \texttt{curated}.
  \item \textbf{Push Docker images} to HCS SWR (Software Repository).
  \item \textbf{Launch ECS instances} (or CCE pods) for api-gateway, ai-service, and pipeline-worker.
  \item \textbf{Configure ELB} with HTTPS listener pointing to api-gateway.
  \item \textbf{Apply security groups}: only ELB accepts external traffic; ECS-to-RDS/OBS is internal only.
  \item \textbf{Run pipeline-worker} to execute the first pipeline on HCS.
  \item \textbf{Validate} by querying the API through the ELB endpoint.
\end{enumerate}

\section{Security Considerations}

\begin{itemize}
  \item \textbf{Network isolation:} VPC with security groups --- no direct internet access to ECS instances.
  \item \textbf{Secrets management:} Database credentials stored in HCS Secret Manager, injected as environment variables.
  \item \textbf{HTTPS termination:} At the ELB layer --- all external traffic is encrypted.
  \item \textbf{Data at rest:} OBS and RDS support encryption at rest (AES-256).
  \item \textbf{Access control:} IAM roles for OBS access; PostgreSQL role-based access for RDS.
\end{itemize}

\section{Benefits of HCS Deployment}

\begin{itemize}
  \item \textbf{Managed services:} RDS handles backups, patching, and high availability. OBS handles replication and durability.
  \item \textbf{Elastic scaling:} ECS instances can scale horizontally for API Gateway under load.
  \item \textbf{Operator co-location:} HCS is deployed on-premises at operator sites --- data never leaves the operator's premises.
  \item \textbf{Compliance:} Tunisian data sovereignty requirements are met by on-premises deployment.
\end{itemize}

\section{Chapter Summary}

This chapter described the cloud deployment strategy: a direct mapping from Docker Compose to HCS services (ECS, OBS, RDS), requiring only environment variable changes. The architecture is cloud-portable by design, with security, scalability, and data sovereignty built in. The next chapter concludes the report.
